\begin{abstract}
We measure how uniformly small language models use their representation space when reading 22 scheduled Indian languages and English. For every layer we compute IsoScore, average random cosine similarity, maximum explainable variance, and intrinsic dimension on token and sentence representations of the n-way parallel IN22-Gen benchmark. We compare two contrastively trained embedding models (EmbeddingGemma-300M, Qwen3-Embedding-0.6B) with two generative models (Gemma-3-1B-PT, Qwen3.5-0.8B-Base). The last transformer layer is near rank-one in most models because of a few massive activation dimensions; the final RMSNorm removes them, so isotropy must be reported before and after it. After the norm, the embedding models are more isotropic than the generative models (token IsoScore 0.23 and 0.11 against 0.04). Languages with high tokenizer fertility, such as Manipuri in Meitei script and Santali in Ol Chiki script, produce the most clustered sentence representations.
\end{abstract}
